% 张三慧《大学物理学》第4版：力学、热学，第1章质点运动学
% 题目唯一源（本候选批次）。PDF 扫描页 53--54，原书印刷页 37--38。
% 保持原书 1.1--1.10 编号；本文件是 LaTeX 片段，不包含 documentclass。
% 转载与正式组卷资格尚待核验，不能自动并入 bank/。
\section*{第1章\quad 质点运动学}
\subsection*{习题（1.1--1.10）}

\paragraph{习题 1.1}
木星的一个卫星——木卫1——上面的塔珥火山喷发出的岩块上升高度可达
\(200\,\mathrm{km}\)。这些石块的喷出速度是多大？
已知木卫1上的重力加速度为 \(1.80\,\mathrm{m/s^2}\)，而且在木卫1上没有空气。

\paragraph{习题 1.2}
一种喷气推进的实验车，从静止开始可在 \(1.80\,\mathrm{s}\) 内加速到
\(1600\,\mathrm{km/h}\) 的速率。按匀加速运动计算，
它的加速度是否超过了人可以忍受的加速度 \(25g\)？
这 \(1.80\,\mathrm{s}\) 内该车跑了多长距离？

\paragraph{习题 1.3}
一辆卡车为了超车，以 \(90\,\mathrm{km/h}\) 的速度驶入左侧逆行道时，
猛然发现前方 \(80\,\mathrm{m}\) 处一辆汽车正迎面驶来。
假定该汽车以 \(65\,\mathrm{km/h}\) 的速度行驶，同时也发现了卡车超车。
设两司机的反应时间都是 \(0.70\,\mathrm{s}\)
（即司机发现险情到实际起动刹车所经过的时间），
他们刹车后的减速度都是 \(7.5\,\mathrm{m/s^2}\)。
试问两车是否会相撞？如果相撞，相撞时卡车的速度多大？

\paragraph{习题 1.4}
跳伞运动员从 \(1200\,\mathrm{m}\) 高空下跳，起初不打开降落伞作加速运动。
由于空气阻力的作用，会加速到一“终极速率” \(200\,\mathrm{km/h}\)
而开始匀速下降。下降到离地面 \(50\,\mathrm{m}\) 处时打开降落伞，
很快速率变为 \(18\,\mathrm{km/h}\) 而匀速下降着地。
若起初加速运动阶段的平均加速度按 \(g/2\) 计，
此跳伞运动员在空中一共经历了多长时间？

\paragraph{习题 1.5}
由消防水龙带的喷嘴喷出的水的流量是 \(q=280\,\mathrm{L/min}\)，
水的流速 \(v=26\,\mathrm{m/s}\)。
若这喷嘴竖直向上喷水，水流上升的高度是多少？
在任一瞬间空中有多少升水？

\paragraph{习题 1.6}
在以初速率 \(v_0=15.0\,\mathrm{m/s}\) 竖直向上抛一块石头后，
\begin{enumerate}
\item 在 \(\Delta t_1=1.0\,\mathrm{s}\) 末又竖直向上抛出第二块石头，
后者在 \(h=11.0\,\mathrm{m}\) 高度处击中前者，求第二块石头抛出时的速率；
\item 若在 \(\Delta t_2=1.3\,\mathrm{s}\) 末竖直向上抛出第二块石头，
它仍在 \(h=11.0\,\mathrm{m}\) 高度处击中前者，
求这一次第二块石头抛出时的速率。
\end{enumerate}

\paragraph{习题 1.7}
一质点在 \(xy\) 平面上运动，运动函数为 \(x=2t,\ y=4t^2-8\)
（采用国际单位制）。
\begin{enumerate}
\item 求质点运动的轨道方程并画出轨道曲线；
\item 求 \(t_1=1\,\mathrm{s}\) 和 \(t_2=2\,\mathrm{s}\) 时，
质点的位置、速度和加速度。
\end{enumerate}

\paragraph{习题 1.8}
男子排球的球网高度为 \(2.43\,\mathrm{m}\)，球网两侧的场地大小都是
\(9.0\,\mathrm{m}\times9.0\,\mathrm{m}\)。
一运动员采用跳发球姿势，其击球点高度为 \(3.5\,\mathrm{m}\)，
离网的水平距离是 \(8.5\,\mathrm{m}\)。
\begin{enumerate}
\item 球以多大速度沿水平方向被击出时，才能使球正好落在对方后方边线上？
\item 球以此速度被击出后过网时，超过网高多少米？
\item 这样，球落地时速率多大？（忽略空气阻力）
\end{enumerate}

\paragraph{习题 1.9}
滑雪运动员离开水平滑雪道飞入空中时的速率
\(v=110\,\mathrm{km/h}\)，着陆的斜坡与水平面夹角
\(\theta=45^\circ\)（见原书图 1.25）。
\begin{enumerate}
\item 计算滑雪运动员着陆时沿斜坡的位移 \(L\) 是多大？
（忽略起飞点到斜面的距离。）
\item 在实际的跳跃中，滑雪运动员所达到的距离
\(L=165\,\mathrm{m}\)，这个结果为什么与计算结果不符？
\end{enumerate}

\paragraph{习题 1.10}
一个人扔石头的最大出手速率 \(v=25\,\mathrm{m/s}\)，
他能把石头扔过与他的水平距离 \(L=50\,\mathrm{m}\)、
高 \(h=13\,\mathrm{m}\) 的一座墙吗？
在这个距离内他能把石头扔过墙的最高高度是多少？


% 第二批：逐页从教材PDF物理页54--55录入，原书印刷页38--39。
% 本批补齐第1章章末习题1.11--1.26；配套解答PDF约20--25页。
\subsection*{续录习题（1.11--1.26）}

\paragraph{习题 1.11}
1997年6月1日，为迎接春节，柯驾车飞越黄河壶口（图1.26）。
东岸跑道长 \(265\,\mathrm{m}\)，柯驾车从跑道东端起动，
到达跑道终端时速度为 \(150\,\mathrm{km/h}\)；
他随即以仰角 \(5^\circ\) 冲出，飞越跨度为
\(57\,\mathrm{m}\)，安全落到西岸木桥上。
\begin{enumerate}
\item 按匀加速运动计算，柯在东岸驾车的加速度和时间各是多少？
\item 柯跨越黄河用了多长时间？
\item 若起飞点高出河面 \(10.0\,\mathrm{m}\)，
柯驾车飞行的最高点离河面几米？
\item 西岸木桥桥面和起飞点的高度差是多少？
\end{enumerate}

\paragraph{习题 1.12}
山上和山下两炮各瞄准对方同时以相同初速各发射一枚炮弹
（原书图1.27），这两枚炮弹会不会在空中相碰？为什么？
忽略空气阻力。如果山高 \(h=50\,\mathrm{m}\)，
两炮相隔的水平距离 \(s=200\,\mathrm{m}\)，
要使这两枚炮弹在空中相碰，它们的速率至少应等于多少？

\paragraph{习题 1.13}
在生物物理实验中用来分离不同种类分子的超级离心机，
转速是 \(6\times10^4\,\mathrm{r/min}\)。
在这种离心机的转子内，离轴 \(10\,\mathrm{cm}\) 远的一个大分子的
向心加速度是重力加速度的几倍？

\paragraph{习题 1.14}
北京天安门所处纬度为 \(39.9^\circ\)，
求它随地球自转的速度和加速度的大小。

\paragraph{习题 1.15}
按玻尔模型，氢原子处于基态时，
它的电子围绕原子核作圆周运动。
电子的速率为 \(2.2\times10^6\,\mathrm{m/s}\)，
离核距离为 \(0.53\times10^{-10}\,\mathrm{m}\)。
求电子绕核运动的频率和向心加速度。

\paragraph{习题 1.16}
北京正负电子对撞机的储存环的周长为 \(240\,\mathrm{m}\)，
电子要沿环以非常接近光速的速率运行。
这些电子运动的向心加速度是重力加速度的几倍？

\paragraph{习题 1.17}
汽车在半径 \(R=400\,\mathrm{m}\) 的圆弧弯道上减速行驶。
设在某一时刻，汽车的速率为 \(v=10\,\mathrm{m/s}\)，
切向加速度的大小为 \(a_t=0.20\,\mathrm{m/s^2}\)。
求汽车的法向加速度和总加速度的大小和方向。

\paragraph{习题 1.18}
一张致密光盘（CD）音轨区域的内半径为
\(R_1=2.2\,\mathrm{cm}\)，外半径为 \(R_2=5.6\,\mathrm{cm}\)
（原书图1.28），径向音轨密度为 \(N=650\,\text{条/mm}\)。
在CD唱机内，光盘每转一圈，激光头沿径向向外移动一条音轨；
激光束相对光盘是以 \(v=1.3\,\mathrm{m/s}\) 的恒定线速度运动的。
\begin{enumerate}
\item 这张光盘的全部放音时间是多少？
\item 激光束到达离盘心 \(r=5.0\,\mathrm{cm}\) 处时，
光盘转动的角速度和角加速度各是多少？
\end{enumerate}

\paragraph{习题 1.19}
一人自由泳时，右手从前到后一次对身体划过的距离
\(\Delta s_{\mathrm{hb}}=1.20\,\mathrm{m}\)，
同时他的身体在泳道中前进了
\(\Delta s_{\mathrm{bw}}=0.9\,\mathrm{m}\)。
求同一时间他的右手在水中划过的距离
\(\Delta s_{\mathrm{hw}}\)。
手对水是向前还是向后划了？

\paragraph{习题 1.20}
当速率为 \(30\,\mathrm{m/s}\) 的西风正吹时，
相对于地面，向东、向西和向北传播的声音的速率各是多大？
已知声音在空气中传播的速率为 \(344\,\mathrm{m/s}\)。

\paragraph{习题 1.21}
一电梯以 \(1.2\,\mathrm{m/s^2}\) 的加速度下降，
其中一乘客在电梯开始下降后 \(0.5\,\mathrm{s}\) 时，
用手在离电梯底板 \(1.5\,\mathrm{m}\) 高处释放一小球。
求此小球落到底板上所需的时间和它对地面下落的距离。

\paragraph{习题 1.22}
一个人骑车以 \(18\,\mathrm{km/h}\) 的速率自东向西行进时，
看见雨点垂直下落；当他的速率增至 \(36\,\mathrm{km/h}\) 时，
看见雨点与他前进的方向成 \(120^\circ\) 下落，
求雨点对地的速度。

\paragraph{习题 1.23}
飞机 \(A\) 以 \(v_A=1000\,\mathrm{km/h}\) 的速率
（相对地面）向南飞行，
同时另一架飞机 \(B\) 以 \(v_B=800\,\mathrm{km/h}\) 的速率
（相对地面）向东偏南 \(30^\circ\) 方向飞行。
求 \(A\) 机相对于 \(B\) 机的速度与 \(B\) 机相对于 \(A\) 机的速度。

\paragraph{习题 1.24}
利用本书中数值表提供的有关数据，
计算原书图1.24中地球表面的大楼日夜相对于太阳参考系的速率之差。
\textit{图示依赖：该题必须结合图1.24所示位置关系；
若未标明纬度，解答需说明采用赤道附近理想化模型。}

\paragraph{习题 1.25}
1964年曾有人做过这样的实验：
测量从以 \(0.99975c\)（\(c\) 为光在真空中的速率，
\(c=2.9979\times10^8\,\mathrm{m/s}\)）的速率运动的
\(\pi^0\) 介子向正前方和正后方发出的光的速率，
测量结果是二者都是 \(c\)。
如果按伽利略变换公式计算，相对于 \(\pi^0\) 介子，
它发出的向正前方和正后方的光的速率应各是多少？

\paragraph{习题 1.26}
曾有报道，当年美国曾用预警飞机帮助以色列的
“爱国者”导弹系统防止伊拉克导弹袭击。
一架预警飞机正在伊拉克上空的速率为
\(150\,\mathrm{km/h}\) 的西风中水平巡航，
机头指向正北，相对于空气的航速为
\(750\,\mathrm{km/h}\)。
飞机中雷达员发现一导弹正相对于飞机以
向西偏南 \(19.5^\circ\) 的方向
以 \(5750\,\mathrm{km/h}\) 的速率水平飞行。
求该导弹相对于地面的速率和方向。
